\usepackage[scr=rsfs,cal=euler]{mathalfa}
\multlinegap0pt
\newcommand{\llbracket}{[\![}
\newcommand{\rrbracket}{]\!]}
\newcommand{\sbullet}{{\scriptscriptstyle\bullet}}
\newcommand{\csbullet}{{\raisebox{1pt}{$\sbullet$}}}
\newcommand\mto{\mathchoice{\longmapsto}{\mapsto}{\mapsto}{\mapsto}}
\newcommand\hto{\mathrel{\lhook\joinrel\to}}
\newcommand\To[1]{\mathchoice{\xrightarrow{\kern1pt#1\kern1pt}}{\stackrel{#1}{\longrightarrow}}{}{}}
\newcommand{\quand}{\quad\text{and}\quad}
\let\epsilon\varepsilon

\usepackage{tikz-cd}
\usepackage[all,cmtip]{xy}
\newdir{ >}{{}*!/-5pt/\dir{>}}

\theoremstyle{plain}
\newtheorem{theorem}{Theorem}[section]
\newtheorem{proposition-definition}{Proposition/Definition}[section]
\newtheorem{theorem-definition}{Theorem/Definition}[section]
\newtheorem*{maintheorem}{Main Theorem}
\newtheorem*{proposition-definition-intro}{Proposition/Definition}
\newtheorem*{definition-intro}{Definition}
\newtheorem*{cor-intro}{Corollary}
\newtheorem*{prop-intro}{Proposition}
\newtheorem{proposition}[theorem]{Proposition}		
\newtheorem{corollary}[theorem]{Corollary}
\newtheorem{lemma}[theorem]{Lemma}
\newtheorem*{conjecture}{Conjecture}
\newtheorem*{conjecture-intro}{Conjecture}
\newtheorem{question}[theorem]{Question}
\newtheorem{claim}[theorem]{Claim}
\newtheorem*{theoremMeta}{Main theorem}
\newtheorem{thm}{Theorem}
\renewcommand*{\thethm}{\Alph{thm}}
\newtheorem{prop}[thm]{Proposition}

\theoremstyle{definition}
\newtheorem{definition}[theorem]{Definition}
\newtheorem{construction-definition}[theorem]{Construction/Definition}
\newtheorem{notation}[theorem]{Notation}
\newtheorem*{problem2prime}{Problem 2'}
\newtheorem{defthm}[thm]{Definition / Theorem}
\newtheorem{problem-Deligne}{Problem}
\newtheorem{remark}[theorem]{Remark}
\newtheorem*{remark-intro}{Remark}
\newtheorem{example}[theorem]{Example}
\newtheorem*{construction}{Construction}


\newcommand{\Abold}{\mathbf{A}}
\newcommand{\Bbold}{\mathbf{B}}
\newcommand{\Cbold}{\mathbf{C}}
\newcommand{\Dbold}{\mathbf{D}}
\newcommand{\Ebold}{\mathbf{E}}
\newcommand{\Kbold}{\mathbf{K}}
\newcommand{\Pbold}{\mathbf{P}}

\newcommand{\NBbb}{\mathbb N}
\newcommand{\PBbb}{\mathbb P}
\newcommand{\QBbb}{\mathbb Q}
\newcommand{\ZBbb}{\mathbb Z}

\newcommand{\Acal}{\mathcal A}
\newcommand{\Bcal}{\mathcal B}
\newcommand{\Ccal}{\mathcal C}
\newcommand{\Dcal}{\mathcal D}
\newcommand{\Ecal}{\mathcal E}
\newcommand{\Gcal}{\mathcal G}
\newcommand{\Ical}{\mathcal I}
\newcommand{\Ocal}{\mathcal O}
\newcommand{\Pcal}{\mathcal P}
\newcommand{\Rcal}{\mathcal R}
\newcommand{\Scal}{\mathcal S}
\newcommand{\Ucal}{\mathcal U}
\newcommand{\Vcal}{\mathcal V}
\newcommand{\Wcal}{\mathcal W}

\newcommand{\Cfrak}{\mathfrak C}
\newcommand{\Hfrak}{\mathfrak H}

\newcommand{\cfrak}{\mathfrak c}
\newcommand{\chfrak}{\mathfrak {ch}}
\newcommand{\sfrak}{\mathfrak s}
\newcommand{\tdfrak}{\mathfrak {td}}

\newcommand{\bis}{\prime\prime}
\newcommand{\tris}{\prime\prime\prime}
\newcommand{\fppf}{\emph{fppf}\xspace}

\newcommand{\Pic}{\mathrm{Pic}}
\newcommand{\prim}{\mathrm{prim}}

\DeclareMathOperator{\Aut}{Aut}
\DeclareMathOperator{\ch}{ch}
\DeclareMathOperator{\CH}{CH}
\DeclareMathOperator{\CHfrak}{\mathfrak{CH}}
\DeclareMathOperator{\Gr}{Gr}
\DeclareMathOperator{\Hom}{Hom}
\DeclareMathOperator{\id}{id}
\DeclareMathOperator{\Imag}{Im}
\DeclareMathOperator{\Picfr}{\mathfrak{Pic}}
\DeclareMathOperator{\iso}{iso}
\DeclareMathOperator{\qiso}{quasi-iso}
\DeclareMathOperator{\rk}{rk}
\DeclareMathOperator{\Spec}{Spec}
\DeclareMathOperator{\Sym}{Sym}
\DeclareMathOperator{\td}{td}
\DeclareMathOperator{\Vect}{Vect}

\let\ov\overline
